\documentclass{beamer}
\usepackage[T1]{fontenc}
\usepackage{lmodern}
\usepackage{graphicx, amsmath, amsthm}
\usepackage{dsfont}
\usepackage[english]{babel}
\usepackage{amssymb,marvosym}
\usepackage{ulem}
\usepackage{cancel}
% Theme auswählen
\usetheme{Madrid}

\newcommand{\set}[1]{\left\{#1\right\}}
\newcommand{\dmid}{\mid\mid} 
\newcommand{\concat}[2]{#1 \cdot #2}
\renewcommand{\angle}[1]{\left\langle #1 \right\rangle}

% Pakete
\usepackage[utf8]{inputenc}
\usepackage[T1]{fontenc}

% Titelinformationen
\title{PAL: Program-aided Language Models}
\author{Eric Peter}
\institute{Universität Heidelberg}
\date{\today}

\begin{document}

% Titelfolie
\begin{frame}
    \titlepage
\end{frame}

% Inhaltsverzeichnis
\begin{frame}{Contents}
    \tableofcontents
\end{frame}

% Abschnitt 1
\section{Introduction}

\begin{frame}{Motivation}
    \begin{block}{Je besser der Prompt, desto besser der Output}
        Resultate verschiedener Prompting-Methoden zeigen:\\
        Der LLM zu erklären, 
        wie sie ihre Ergebnisse richtig herleiten soll,
        spielt eine unüberschaubare Rolle für den generierten Output. 
    \end{block}

    % \pause 

    \underline{Bisherige Methoden}:
    \begin{itemize}
        \item \textbf{Few-Shot-Prompting}
        \item \textbf{Chain-Of-Thought-Prompting}
    \end{itemize}

    % \pause

    \underline{Zielsetzung}:\\
    In diesen Thema geht es um eine weitere Prompting-Methode,
    die in Kombination mit bisherigen Methoden
    eine neue \textit{state-of-the-art}-Benchmark setzt.\\

    % \pause

    Es geht um das \begin{center}
        \large \textbf{PAL: Program-aided Language Models}
    \end{center} 
    
\end{frame}

\begin{frame}{Background: Few-Shot-Prompting}

    \begin{block}{Beispiellösungen helfen!}
        \begin{itemize}
            \item Füttere die Maschine mit \(k\) beispielhaften Frage-Antwort-Paaren \(\set{(x_i,y_i)}_{i=1}^k\).
            \item Gestalte die Demonstrationen in einem einheitlichen Format,
            sodass die LLM das leicht das Muster erkennt, 
            dem sie folgen soll.
            \item Die Beispiele stehen im Kontext zur jeweiligen Aufgabe und zeigen, welche verallgemeinernden Fähigkeiten oder Kenntnisse nötig sind.
        \end{itemize}
    \end{block}

    Wir konstruieren einen Prompt\[
        p \equiv \angle{\concat{x_1}{y_1}} \dmid \angle{\concat{x_2}{y_2}} \dmid ... \dmid \angle{\concat{x_k}{y_k}},
    \]wobei Frage \(x_i\) und Lösung \(y_i\) mit dem Zeichen '\(\cdot\)' 
    und die Beispiele mit '\(\dmid\)' geeignet konkateniert werden.\\
    Anschließend fügen wir die Testinstanz \(x_{test}\) hinzu 
    und lassen die Maschine mit \(p \dmid x_{test}\) den Output \(y_{test}\) generieren.

\end{frame}

\begin{frame}{Background: Chain-Of-Thought-Prompting}
    \begin{block}{Erweitere die Parameter des Few-Shot-Promptings}
        Neben den bisherigen Parametern \(x_i\) und \(y_i\) 
        wollen wir durch Einfügen geeigneter \(t_i\) der Maschine beibringen, \begin{itemize}
            \item das Problem in weniger komplexe Teilprobleme zu teilen,
            \item Teilprobleme einzeln zu lösen und auf Basis der Zwischenergebnisse weiterzuarbeiten
            \item und diese Unterteilung im Output darzustellen.
        \end{itemize}
    \end{block}

    Konkret erstellen wir \[
        p \equiv \angle{\concat{\concat{x_1}{t_1}}{y_1}} \dmid \angle{\concat{\concat{x_2}{t_2}}{y_2}} \dmid ... \dmid \angle{\concat{\concat{x_k}{t_k}}{y_k}}
    \]und lassen nun mit \(p \dmid x_{test}\) sowohl die Erklärung \(t_{test}\) als auch die Lösung \(y_{test}\) generieren.
\end{frame}

\begin{frame}{Zielsetzung}
    Ziel dieser Seminararbeit ist es,
    ein bestimmtes Thema zu analysieren und zu präsentieren.
\end{frame}

% Abschnitt 2
\section{Hauptteil}

\begin{frame}{Wichtige Ergebnisse}
    \begin{block}{Kernaussage}
        Hier steht eine wichtige Erkenntnis.
    \end{block}

    \pause

    \begin{alertblock}{Achtung}
        Hier kann ein kritischer Punkt hervorgehoben werden.
    \end{alertblock}
\end{frame}

% Abschnitt 3
\section{Fazit}

\begin{frame}{Fazit}
    \begin{enumerate}
        \item Wichtigstes Ergebnis
        \item Erkenntnisse
        \item Ausblick
    \end{enumerate}
\end{frame}

% Letzte Folie
\begin{frame}
    \centering
    \Huge Vielen Dank!
\end{frame}

\end{document}
